\subsection{准素分解与纯量扩张}

本节中，A 和 B 表示两个环；我们考虑环同态\(\rho: \mathbf{A} \to \mathbf{B}\)，它使 B 成为 A-代数；回忆，对于每个 B-模 F，\(\rho_*(\mathbf{F})\)是交换群 F，其 A-模结构由\(a.y = \rho(a)y\)对所有\(a \in \mathbf{A}, y \in \mathbf{F}\)定义。

引理 1. 设 A 是诺特环，p 是 A 的素理想，E 是 A-模，其湮灭子包含 p 的某个幂，且 \(\operatorname{Ass}(\mathbf{E}) = \{\mathfrak{p}\}\)；F 是 B-模，且 \(\rho_{*}(\mathbf{F})\) 是平坦的 A-模。则条件 \(\mathfrak{P} \in \operatorname{Ass}_{\mathbf{B}}(\mathbf{E} \otimes_{\mathbf{A}} \mathbf{F})\) 蕴含 \(\overline{\rho}^{1}(\mathfrak{P}) = \mathfrak{p}\)。

若\(n\)满足\(p^n\mathbf{E} = 0\)，则\(p^n\mathbf{B} \subset \operatorname{Ann}(\mathbf{E} \otimes_{\mathbf{A}} \mathbf{F})\)，从而\(p^n\mathbf{B} \subset \mathfrak{P}\)，这意味着\(p^n \subset \overline{\rho}^1(\mathfrak{P})\)，因而\(p \subset \overline{\rho}^1(\mathfrak{P})\)，因为\(\overline{\rho}^1(\mathfrak{P})\)是素理想。此外，若\(a \in \mathbf{A} - p\)，则\(h\)上的伸缩映射比为\(a\)，且在\(\mathbf{E}\)上是单射（§ 1，第 1 节，命题 2 的推论 2）；由于\(h \otimes l_{\mathbf{F}}\)是作用于\(h'\)的伸缩映射，其比为\(\rho(a)\)，作用于\(\mathbf{E} \otimes_{\mathbf{A}} \mathbf{F}\)上，且\(\rho_*(\mathbf{F})\)是平坦的，故\(h'\)是单射（第一章，第 2 节，第 2 节，定义 1）；这证明\(\rho(a) \notin \mathfrak{P}\)，从而\(\overline{\rho}^1(\mathfrak{P}) = p\)。

定理 2. 设 \(\rho: \mathbf{A} \to \mathbf{B}\) 是环同态，E 是 A-模，F 是 B-模，且 \(\rho_*(\mathbf{F})\) 是平坦的 A-模。则

\[
\operatorname{Ass} _ {\mathbf {B}} (\mathrm{E} \otimes_ {\mathbf {A}} \mathrm{F}) \supset \bigcup_ {\mathfrak {p} \in \operatorname{Ass} _ {\mathbf {A}} (\mathrm{E})} \operatorname{Ass} _ {\mathbf {B}} (\mathrm{F} / \mathfrak {p} \mathrm{F}).\tag{5}
\]

当 A 是诺特环时，(5) 两边相等。

令 \(\mathfrak{p} \in \operatorname{Ass}_{\mathbf{A}}(\mathbf{E})\)；根据定义，存在正合列

\[
0 \rightarrow \mathrm{A} / \mathfrak {p} \rightarrow \mathrm{E}.
\]

由于 F 是平坦的 A-模，我们得到正合列

\[
0 \rightarrow \mathrm{F} / \mathfrak {p} \mathrm{F} \rightarrow \mathrm{E} \otimes_ {\mathrm{A}} \mathrm{F}
\]

从而 \(\operatorname{Ass}_{\mathbf{B}}(\mathbf{F} / p\mathbf{F})\subset \operatorname{Ass}_{\mathbf{B}}(\mathbf{E}\otimes_{\mathbf{A}}\mathbf{F})\)，这就证明了包含关系 (5)。

现在假设 A 是诺特环，我们来证明反向包含关系。

我们分步进行：

\begin{enumerate}
\def\labelenumi{(\roman{enumi})}
\tightlist
\item
  首先假设 E 是有限生成的 A-模，且 \(\operatorname{Ass}_{\mathbf{A}}(\mathbf{E})\) 归结为单个元素 \(p\)。根据 § 1，第 4 节，定理 1，存在 E 的合成列 \((\mathrm{E}_i)_{0 \leqslant i \leqslant n}\)，使得 \(\mathrm{E}_i / \mathrm{E}_{i+1}\) 同构于 \(\mathbf{A} / p_i\)，其中 \(p_i\) 是 A 的素理想；此外（§ 1，第 4 节，定理 2 以及第 3 节命题 7），所有 \(p_i\) 都包含 \(p\)。由于 F 是平坦的 A-模，\(\mathrm{E}_i \otimes_{\mathbf{A}} \mathbf{F}\) 构成 \(\mathbf{E} \otimes_{\mathbf{A}} \mathbf{F}\) 的合成列，且 \((\mathrm{E}_i \otimes_{\mathbf{A}} \mathbf{F}) / (\mathrm{E}_{i+1} \otimes_{\mathbf{A}} \mathbf{F})\) 可识别为
\end{enumerate}

\[
(\mathbf {A} / \mathfrak {p} _ {i}) \otimes_ {\mathbf {A}} \mathbf {F} = \mathbf {F} / \mathfrak {p} _ {i} \mathbf {F}.
\]

于是根据 § 1，第 1 节，命题 3，

\[
\operatorname{Ass} _ {\mathbf {B}} (\mathrm{E} \otimes_ {\mathbf {A}} \mathrm{F}) \subset \bigcup_ {i = 0} ^ {n - 1} \operatorname{Ass} _ {\mathbf {B}} (\mathrm{F} / p _ {i} \mathrm{F}).
\]

  我们知道 E 被 p 的某个幂湮灭（第 1 节，注）；于是引理 1 表明，对于所有 \(\mathfrak{P} \in \operatorname{Ass}_{\mathbf{B}}(\mathbf{E} \otimes_{\mathbf{A}} \mathbf{F})\)，都有 \(\overline{\rho}^{1}(\mathfrak{P}) = \mathfrak{p}\)。由于 \(\mathbf{F} / \mathfrak{p}_i \mathbf{F}\) 同构于 \((\mathbf{A} / \mathfrak{p}_i) \otimes_{\mathbf{A}} \mathbf{F}\)，根据引理 1，\(\overline{\rho}^{1}(\mathfrak{P}') = \mathfrak{p}_i\) 对所有 \(\mathfrak{P}' \in \operatorname{Ass}_{\mathbf{B}}(\mathbf{F} / \mathfrak{p}_i \mathbf{F})\) 都成立；因此，\(\operatorname{Ass}_{\mathbf{B}}(\mathbf{E} \otimes_{\mathbf{A}} \mathbf{F}) \cap \operatorname{Ass}(\mathbf{F} / \mathfrak{p}_i \mathbf{F}) = \varnothing\) 当 \(\mathfrak{p}_i \neq \mathfrak{p}\) 时成立，从而当前情形下定理得证。

\begin{enumerate}
\def\labelenumi{(\roman{enumi})}
\setcounter{enumi}{1}
\tightlist
\item
  只假设 E 是有限生成的 A-模。令 \(p_{i}\)（\(1 \leqslant i \leqslant n\)）为 \(\operatorname{Ass}_{\mathbf{A}}(\mathbf{E})\) 的元素，并令 \(\{0\} = \bigcap_{i=1}^{n} Q_{i}\) 为相应的既约准素分解（第 3 节）；于是 E 同构于直和中 \(E_{i} = E/Q_{i}\) 的一个子模，而由于 F 是平坦的 A-模，\(E \otimes_{A} F\) 同构于 B-模 \(E_{i} \otimes_{A} F\) 的直和的一个子模。由此（§ 1，第 1 节，命题 3 及其推论 1），我们得到
\end{enumerate}

\[
\operatorname{Ass} _ {\mathbf {B}} (\mathrm{E} \otimes_ {\mathrm{A}} \mathrm{F}) \subset \bigcup_ {i = 1} ^ {n} \operatorname{Ass} _ {\mathbf {B}} (\mathrm{E} _ {i} \otimes_ {\mathrm{A}} \mathrm{F}).
\]

但是 \(E_{i}\) 是有限生成的 A-模，且 \(\operatorname{Ass}_{\mathbf{A}}(\mathbf{E}_{i})\) 归结为单个元素 \(p_{i}\)（第 1 节，定义 1）。根据 (i)，\(\operatorname{Ass}_{\mathbf{B}}(\mathbf{E}_{i} \otimes_{\mathbf{A}} \mathbf{F}) = \operatorname{Ass}_{\mathbf{B}}(\mathbf{F}/p_{i}\mathbf{F})\)，定理在此情形得证。

\begin{enumerate}
\def\labelenumi{(\roman{enumi})}
\setcounter{enumi}{2}
\tightlist
\item
  一般情形。B-模 \(E \otimes_{A} F\) 是子模 \(E' \otimes_{A} F\) 的并，其中 \(E'\) 遍历 A-模 E 的有限生成子模所成的集合。若 P 属于 \(\operatorname{Ass}_{\mathbf{B}}(\mathbf{E} \otimes_{\mathbf{A}} \mathbf{F})\)，则存在 E 的有限生成子模 \(E'\)，使得 \(\mathfrak{P} \in \operatorname{Ass}_{\mathbf{B}}(\mathbf{E}' \otimes_{\mathbf{A}} \mathbf{F})\)。根据 (ii)，存在 \(\mathfrak{p} \in \operatorname{Ass}_{\mathbf{A}}(\mathbf{E}')\)，使得 \(\mathfrak{P} \in \operatorname{Ass}_{\mathbf{B}}(\mathbf{F}/\mathfrak{p}\mathbf{F})\)；由于 \(\operatorname{Ass}_{\mathbf{A}}(\mathbf{E}') \subset \operatorname{Ass}_{\mathbf{A}}(\mathbf{E})\)，定理 2 的证明完成。
\end{enumerate}

推论 1. 若 A 是诺特环，且 \(\mathfrak{P} \in \operatorname{Ass}_{\mathbf{B}}(\mathbf{E} \otimes_{\mathbf{A}} \mathbf{F})\)，则 \(\overline{\rho}^{1}(\mathfrak{P}) \in \operatorname{Ass}_{\mathbf{A}}(\mathbf{E})\)，并且 \(\overline{\rho}^{1}(\mathfrak{P})\) 是 A 中的唯一素理想 \(\mathfrak{p}\)，满足 \(\mathfrak{P} \in \operatorname{Ass}_{\mathbf{B}}(\mathbf{F}/\mathfrak{p}\mathbf{F})\)。

这由定理 2 以及引理 1 应用于 \(\mathbf{E} = \mathbf{A} / \mathfrak{p}\) 的情形立即得出。

推论 2. 假设 A 和 B 都是诺特环，且 B 是平坦的 A-模。令 p 是 A 的素理想，Q ⊂ E 是 p-准素子模，P 是 B 的素理想。要使 Q ⊗A B 成为 E ⊗A B 的 P-准素子模，充要条件是 pB 是 B 的 P-准素理想。

  将定理 2 应用于 A-模 E/Q 和 B-模 B；则 \(\operatorname{Ass}_{\mathbf{A}}(\mathrm{E} / \mathrm{Q}) = \{\mathfrak{p}\}\)，且 \((\mathrm{E} / \mathrm{Q})\otimes_{\mathbf{A}}\mathbf{B}\) 同构于 \((\mathrm{E}\otimes_{\mathbf{A}}\mathrm{B}) / (\mathrm{Q}\otimes_{\mathbf{A}}\mathrm{B})\)，从而 \(\operatorname{Ass}_{\mathbf{B}}((\mathrm{E}\otimes_{\mathbf{A}}\mathrm{B}) / (\mathrm{Q}\otimes_{\mathbf{A}}\mathrm{B})) = \operatorname{Ass}_{\mathbf{B}}(\mathrm{B} / \mathfrak{p}\mathbf{B})\)。称 \(\mathbf{Q}\otimes_{\mathbf{A}}\mathbf{B}\) 为 \(\mathfrak{P}\)-准素子模（它是 \(\mathrm{E}\otimes_{\mathbf{A}}\mathrm{B}\) 的子模），等价于 \(\operatorname{Ass}_{\mathbf{B}}(\mathrm{B} / \mathfrak{p}\mathrm{B})\) 归结为 \(\mathfrak{P}\)，推论得证。

  注。假设 A 和 B 都是诺特环。设 \(\mathfrak{P}\) 是 B 的素理想，令 \(\mathfrak{p} = \overline{\rho}^{1}(\mathfrak{P})\)；记 \(S = A - \mathfrak{p}\)，并令 \(k(\mathfrak{p}) = S^{-1}(A / \mathfrak{p})\) 为 \(A / \mathfrak{p}\) 的分式域。由于 \(\mathfrak{P}\) 包含 \(\mathfrak{pB}\)，\(\mathfrak{P} / \mathfrak{pB}\) 是 \(B / \mathfrak{pB}\) 的素理想。若 \(\rho'\) 是复合同态 \(A \xrightarrow{\rho} B \to B / \mathfrak{pB}\)，则我们知道 \(S^{-1}(B / \mathfrak{pB})\) 可识别为环 \((\rho'(S))^{-1}(B / \mathfrak{pB})\)，且 \(\mathfrak{P}' = S^{-1}(\mathfrak{P} / \mathfrak{pB})\) 可识别为该环的一个理想（第二章，第 2 节，第 2 号，命题 6）；由于 \(\mathfrak{P} / \mathfrak{pB}\) 不与 \(\rho'(S)\) 相交，\(\mathfrak{P}'\) 是 \(S^{-1}(B / \mathfrak{pB})\) 的素理想（第二章，第 2 节，第 5 号，命题 11）；此外，\(S^{-1}(B / \mathfrak{pB})\)、\(S^{-1}((A / \mathfrak{p}) \otimes_{A} B)\) 和 \((S^{-1}(A / \mathfrak{p})) \otimes_{A} B = k(\mathfrak{p}) \otimes_{A} B\) 之间有典范同构；类似地，\(S^{-1}(F / \mathfrak{pF})\) 可典范地识别为 \(k(\mathfrak{p}) \otimes_{A} F\)。因此，在定理 2 的假设下，要使 \(\mathfrak{P} \in \operatorname{Ass}_{B}(E \otimes_{A} F)\)，充要条件是 \(\mathfrak{p} \in \operatorname{Ass}_{A}(E)\) 以及

\[
\mathfrak {P} ^ {\prime} \in \operatorname{Ass} _ {k (\mathfrak {p}) \otimes_ {\mathbf {A}} \mathbf {B}} (k (\mathfrak {p}) \otimes_ {\mathbf {A}} \mathbf {F}).
\]

因为根据定理 2 及其推论 1，这归结为验证条件

\[
" \mathfrak {P} \in \operatorname{Ass} _ {\mathbf {B}} (\mathrm{F} / \mathfrak {p} \mathrm{F})" \quad \text { and } \quad " \mathfrak {P} ^ {\prime} \in \operatorname{Ass} _ {k (\mathfrak {p}) \otimes_ {\mathbf {A}} \mathbf {B}} (k (\mathfrak {p}) \otimes_ {\mathbf {A}} \mathrm{F})"
\]

这些条件是等价的；但由于 B 是诺特环，这由 § 1，第 2 节，命题 5 的推论以及上述识别立即得出。

命题 11. 假设 A 和 B 都是诺特环，且 B 是平坦的 A-模。设 E 是 A-模，E' 是 E 的子模，且对于每个理想 \(\mathfrak{p} \in \operatorname{Ass}_{\mathbf{A}}(\mathbf{E}/\mathbf{E}')\)，\(\mathfrak{p}\mathbf{B}\) 是 B 的素理想或等于 B。设 \(\mathbf{E}' = \bigcap_{\mathfrak{p} \in \operatorname{Ass}(\mathbf{E}/\mathbf{E}')} \mathbf{Q}(\mathfrak{p})\) 是 \(\mathbf{E}'\) 在 E 中的既约准素分解，其中对于所有 \(\mathbf{Q}(\mathfrak{p})\)，\(\mathfrak{p}\) 都是 \(\mathfrak{p} \in \operatorname{Ass}(\mathbf{E}/\mathbf{E}')\)-准素的。

\begin{enumerate}
\def\labelenumi{(\roman{enumi})}
\item
  若 \(\mathfrak{p} \in \operatorname{Ass}(\mathrm{E} / \mathrm{E}')\) 且 \(\mathfrak{p}\mathbf{B} = \mathbf{B}\)，则 \(\mathbf{Q}(\mathfrak{p}) \otimes_{\mathbf{A}} \mathbf{B} = \mathbf{E} \otimes_{\mathbf{A}} \mathbf{B}\)。
\item
  若 \(\mathfrak{p} \in \operatorname{Ass}(\mathrm{E} / \mathrm{E}')\) 且 \(\mathfrak{pB}\) 是素理想，则 \(\mathsf{Q}(\mathfrak{p})\) 在 \(\mathfrak{pB}\) 中是 \(\mathbf{E} \otimes_{\mathbf{A}} \mathbf{B}\)-准素的。
\item
  若 \(\Phi\) 是满足 \(\mathfrak{p} \in \operatorname{Ass}(\mathrm{E} / \mathrm{E}')\) 为素理想的 \(\mathfrak{pB}\) 所成的集合，则
\end{enumerate}

\[
\mathbf {E} ^ {\prime} \otimes_ {\mathbf {A}} \mathbf {B} = \bigcap_ {\mathfrak {p} \in \Phi} (\mathbf {Q} (\mathfrak {p}) \otimes_ {\mathbf {A}} \mathbf {B})
\]

并且这个关系是 \(E^{\prime} \otimes_{A} B\) 在 \(E \otimes_{A} B\) 中的既约准素分解。

若 \(\mathfrak{pB} = \mathbf{B}\)，将定理 2 应用于 \(\mathbf{E} / \mathbf{Q}(\mathfrak{p})\) 和 \(\mathbf{B}\)，得到

\[
\operatorname{Ass} _ {\mathbf {B}} ((\mathrm{E} / \mathrm{Q} (\mathfrak {p})) \otimes_ {\mathbf {A}} \mathbf {B}) = \varnothing
\]

并且，由于 B 是诺特环且是平坦的 A-模，我们得到（第 1 节，第 1 号，命题 2 的推论 1）\(\mathbf{Q}(\mathfrak{p}) \otimes_{\mathbf{A}} \mathbf{B} = \mathbf{E} \otimes_{\mathbf{A}} \mathbf{B}\)。断言 (ii) 由定理 2 的推论 2 得出，取 \(\mathfrak{P} = \mathfrak{p}\mathbf{B}\)。最后，关系 \(\mathbf{E}' \otimes_{\mathbf{A}} \mathbf{B} = \bigcap_{\mathfrak{p} \in \Phi} (\mathbf{Q}(\mathfrak{p}) \otimes_{\mathbf{A}} \mathbf{B})\) 由 B 是平坦的 A-模这一事实得出（第一章，第 2 节，第 6 号，命题 6）；由于 \(\mathfrak{p} = \overline{\rho}^{1}(\mathfrak{p}\mathbf{B})\) 对于 \(\mathfrak{p} \in \Phi\) 成立（引理 1），有 \(\mathfrak{p}\mathbf{B} \neq \mathfrak{p}'\mathbf{B}\)，其中两个不同的理想为 \(\mathfrak{p}, \mathfrak{p}'\)，且属于集合 \(\Phi\)；另一方面，

\[
\operatorname{Ass} ((\mathbf {E} \otimes_ {\mathbf {A}} \mathbf {B}) / (\mathbf {E} ^ {\prime} \otimes_ {\mathbf {A}} \mathbf {B})) = \Phi
\]

是一个既约准素分解。

\[
\mathbf {E} ^ {\prime} \otimes_ {\mathbf {A}} \mathbf {B} = \bigcap_ {\mathfrak {p} \in \Phi} (\mathbf {Q} (\mathfrak {p}) \otimes_ {\mathbf {A}} \mathbf {B})
\]

  是一个既约准素分解。

推论。假设 \(\mathfrak{pB}\) 对所有 \(\mathfrak{p} \in \operatorname{Ass}_{\mathbf{A}}(\mathrm{E}/\mathrm{E}')\) 都是素理想。于是，若 \(\mathfrak{p}_1, \ldots, \mathfrak{p}_n\) 是 \(\operatorname{Ass}_{\mathbf{A}}(\mathrm{E}/\mathrm{E}')\) 的极小元，则 \(\mathfrak{p}_i\mathbf{B}\) 是

\[
\operatorname{Ass} _ {\mathbf {A}} ((\mathbf {E} \otimes_ {\mathbf {A}} \mathbf {B}) / (\mathbf {E} ^ {\prime} \otimes_ {\mathbf {A}} \mathbf {B})).
\]

  由命题 11 可知，在此情形下有 \(\mathfrak{p}_i\mathbf{B} \neq \mathfrak{p}_j\mathbf{B}\)，其中 \(i \neq j\)。

\textbf{例}\par

\begin{enumerate}
\def\labelenumi{(\arabic{enumi})}
\item
  取 \(B = S^{-1}A\)，其中 S 是 A 的乘性子集；若 A 是诺特环，则命题 11 的假设满足，从而得到第 4 节命题 6 的一部分。
\item
  设 A 是诺特环，m 是 A 的一个理想，B 是 A 关于 m-进拓扑的 Hausdorff 完备化；则 B 是平坦的 A-模，可以将定理 2 应用于 F = B；但一般而言，对于 A 的素理想，命题 11 的假设并不满足（第三章，第 2 节，习题 15 (b)）。
\item
  设 A 是诺特环，B 是多项式代数 \(A[X_{1},\ldots,X_{n}]\)；B 是诺特环且是自由 A-模，因而平坦。此外，若 p 是 A 的素理想，则 B/pB 同构于 \((A/p)[X_{1},\ldots,X_{n}]\)，后者是整环，所以 pB 是素理想；因此，命题 11 的假设对每个 A-模 E 及其每个子模 \(E'\) 都满足。
\item
  设 A 是域 k 上的有限生成代数，K 是 k 的扩域，\(B = A \otimes_{k} K\) 是通过纯量扩张得到的 K-代数；A 和 B 都是诺特环，且 B 是自由的 A-模，因此定理 2 可应用于 F = B。在某些情形（例如 k 代数闭时），可以证明 A 的每个素理想 p 都使得 pB 为素理想或等于 B；我们稍后还会回到这个例子。
\end{enumerate}

